\subsection{Постановка задачи и результат работы}
По характеристикам дома необходимо предсказать его стоимость
\file{SalePrice}. Пусть \(x_i\) --- вектор подготовленных признаков,
\(y_i\) --- известная цена, а \(f_\theta\) --- модель с обучаемыми
параметрами \(\theta\). Тогда прогноз имеет вид
\begin{equation}
\widehat y_i=f_\theta(x_i),\qquad y_i\in\mathbb R_{>0}.
\end{equation}
Требуется подготовить данные, построить последовательную нейронную сеть,
исследовать её настройки, оценить результат на выделенной валидационной
части и получить прогнозы для домов из \file{test.csv}.

\textbf{Результат сдачи:} выполненный ноутбук, таблица экспериментов,
диагностические графики, обоснованный выбор конфигурации и таблица
\file{Id, SalePrice}. Отдельно сохраняются исходные данные эксперимента
и сведения, позволяющие его повторить.

\subsection{Обучение, валидация и данные для прогнозирования}
В исходном файле \file{train.csv} известны правильные ответы.
Разделим его строки на обучающую часть (80\%) и валидационную (20\%).
Первую используем для определения параметров преобразований и весов сети,
вторую --- для сравнения конфигураций. В \file{test.csv} ответов нет.

\begin{keybox}{Основное правило}
Медианы, средние, масштабы и словарь категорий определяются по обучающей
части. К валидации и новым домам применяются уже найденные преобразования.
В терминах scikit-learn: \file{fit\_transform} для обучения,
\file{transform} для остальных данных~\cite{pitfalls}.
\end{keybox}

Если сначала вычислить медиану по всему \file{train.csv}, а затем
разделить строки, в подготовку обучения попадёт информация о валидации.
Если отдельно закодировать категории в двух таблицах, один и тот же
числовой код может означать разные категории. Оба случая нарушают
условия корректного сравнения.

Случайное разделение подходит для учебного предположения о независимых
однородных наблюдениях. Для временного прогноза, повторных продаж одного
объекта или групп зависимых домов нужен другой способ разбиения.
Укажите применяемое предположение в отчёте.

Повторный выбор модели по одной валидации адаптирует решение к ней.
Полученные числа являются оценкой на использованном разбиении, а не
независимой финальной проверкой. Для такой проверки понадобилась бы
отдельная размеченная выборка, не участвовавшая в выборе настроек.

\subsection{Пропуски, категории и масштабирование}
Для числового признака в исходном варианте будем заменять пропуски
медианой обучающей части. Медиана меньше реагирует на отдельные очень
большие значения, чем среднее. Выбранное правило одинаково применяется
ко всем частям данных~\cite{imputer}.

Для категорий введём отдельное значение \file{\_\_missing\_\_}.
По описанию данных необходимо проверить, означает ли отсутствие записи
неизвестное значение или отсутствие самого объекта: например, гаража.
Автоматически отождествлять эти случаи нельзя. В исходном варианте
сохраняется фиксированное удаление пяти столбцов из задания:
\file{Alley}, \file{FireplaceQu}, \file{PoolQC}, \file{Fence},
\file{MiscFeature}. Это условие учебного сравнения, а не универсальное
правило обработки пропусков.

Категории без естественного порядка кодируем методом \textit{one-hot}:
каждой известной категории соответствует отдельный бинарный столбец.
Порядок столбцов задаётся кодировщиком, обученным один раз.
\file{handle\_unknown="ignore"} позволяет обработать ранее не встречавшуюся
категорию: её блок индикаторов будет нулевым~\cite{onehot}.
Это не даёт модели знаний о новой категории; её появление стоит отметить.

\begin{keybox}{Небольшой пример}
Для категорий \(A,B,C\) возможны коды
\(A\mapsto(1,0,0)\), \(B\mapsto(0,1,0)\),
\(C\mapsto(0,0,1)\). Замена их числами \(1,2,3\) навязывает
порядок и расстояния, которых в описании признака может не быть.
\end{keybox}

Числовые признаки приведём к сопоставимому масштабу:
\begin{equation}
\widetilde x_{ij}=\frac{x_{ij}-\mu_j}{s_j},
\end{equation}
где \(\mu_j,s_j\) вычислены только по обучающей части после заполнения
пропусков. Для постоянного признака реализация использует безопасный
масштаб 1. Масштабирование не делает признаки независимыми и не
превращает распределение в нормальное~\cite{scaler}.

В примере также стандартизуется целевая цена:
\begin{equation}\label{eq:target}
z_i=\frac{y_i-\mu_y}{s_y},\qquad
\widehat y_i=\mu_y+s_y\widehat z_i.
\end{equation}
Это помогает обучать сеть на числах умеренного масштаба.
Перед оценкой и сохранением прогноз возвращается к исходной цене.
Использовать среднюю цену валидационной части при обратном преобразовании
нельзя.

\subsection{Полносвязная нейронная сеть}
Нейрон полносвязного слоя вычисляет взвешенную сумму и применяет
функцию активации:
\begin{equation}
a_j=\sum_{k=1}^{d}w_{jk}x_k+b_j,\qquad h_j=\varphi(a_j).
\end{equation}
В матричной записи для слоя \(\ell\):
\begin{equation}
h^{(\ell)}=\varphi\!\left(W^{(\ell)}h^{(\ell-1)}+b^{(\ell)}\right),
\qquad h^{(0)}=\widetilde x.
\end{equation}
Сеть с одним скрытым слоем и линейным выходом задаёт
\begin{equation}
\widehat z=v^{\mathsf T}\varphi(W\widetilde x+b)+c.
\end{equation}
Реализация последовательности таких слоёв соответствует
\file{keras.Sequential}; полносвязный слой представлен
\file{Dense}~\cite{sequential,dense}.

\begin{figure}[htbp]
\centering\begin{tikzpicture}[>=Stealth,
  neuron/.style={circle,draw=accent,fill=light,minimum size=9mm},
  every node/.style={font=\small}]
\foreach \i/\label in {1/1,2/j,3/d}{
  \node[neuron] (x\i) at (0,-1.2*\i) {$x_{\label}$};}
\foreach \i/\label in {1/1,2/k,3/m}{
  \node[neuron] (h\i) at (4,-1.2*\i) {$h_{\label}$};}
\node[neuron] (y) at (8,-2.4) {$\widehat z$};
\foreach \i in {1,2,3}{\foreach \j in {1,2,3}{
  \draw[->,black!30] (x\i)--(h\j);}}
\foreach \i in {1,2,3}{\draw[->,black!45] (h\i)--(y);}
\node[align=center] at (0,0) {Подготовленные\\признаки};
\node[align=center] at (4,0) {Скрытый слой\\ReLU};
\node[align=center] at (8,0) {Линейный\\выход};
\end{tikzpicture}

\caption{Схема полносвязной сети. Показаны отдельные компоненты;
фактическое число входов определяется после предобработки.}
\label{fig:network}
\end{figure}

Входная размерность \(d\) равна числу подготовленных признаков.
Её нельзя произвольно менять при фиксированных данных.
В задании об изменении числа нейронов варьируется ширина
\textbf{первого скрытого полносвязного слоя}. Выходной слой содержит
один нейрон, поскольку для каждого дома предсказывается одно число.

У слоя с \(d\) входами и \(m\) нейронами имеется \(dm+m\)
обучаемых параметров. Для схемы \(d\to m\to1\) их всего
\begin{equation}
P=(d+1)m+(m+1).
\end{equation}
Например, при \(d=10\), \(m=8\) получаем \(P=97\).
Это вычислительная иллюстрация, а не характеристика файлов курса.
Сравните собственный расчёт с \file{model.count\_params()}.

\subsection{Функции активации}
Нелинейная активация позволяет сети представлять нелинейные зависимости.
Последовательность только линейных слоёв без активаций сводится к одному
аффинному преобразованию --- взвешенной сумме входов с добавлением смещения.

\noindent\begin{tabularx}{\textwidth}{@{}p{0.16\textwidth}p{0.28\textwidth}Y@{}}
\toprule
Активация & Формула & Особенность \\
\midrule
ReLU & \(\max(0,a)\) & Простая исходная активация скрытых слоёв;
на отрицательной полуоси производная равна нулю. \\
Sigmoid & \(1/(1+e^{-a})\) & Выход в \((0,1)\); в областях насыщения
производная мала. \\
Tanh & \(\tanh(a)\) & Выход в \((-1,1)\); также имеет области насыщения. \\
Линейная & \(a\) & Не ограничивает диапазон прогноза. \\
\bottomrule
\end{tabularx}

В основной конфигурации используем ReLU внутри сети и линейную функцию
на выходе~\cite{activations}. Sigmoid на выходе без согласованного
преобразования цели ограничила бы прогноз интервалом \((0,1)\).
Линейный выход для стандартизованной цели может принимать отрицательные
значения: это означает цену ниже \(\mu_y\), а не обязательно отрицательную
цену. Проверять допустимость нужно после обратного преобразования.

\subsection{Функции потерь и метрики регрессии}
Пусть на оцениваемой части имеется \(n\) объектов, а ошибка прогноза
\(e_i=y_i-\widehat y_i\). Определим
\begin{align}
\operatorname{MSE}&=\frac1n\sum_{i=1}^{n}e_i^2,\\
\operatorname{MAE}&=\frac1n\sum_{i=1}^{n}|e_i|,\\
\operatorname{RMSE}&=\sqrt{\frac1n\sum_{i=1}^{n}e_i^2}.
\end{align}
MAE и RMSE измеряются в единицах цены, MSE --- в квадрате этих единиц.
Квадратичная ошибка сильнее выделяет большие промахи. MAE придаёт
линейный вес абсолютной величине ошибки~\cite{losses}.

\begin{keybox}{Пример расчёта}
Пусть цены равны \((100,150,200)\), а прогнозы --- \((110,140,230)\)
в условных денежных единицах. Тогда
\(\operatorname{MAE}=50/3\approx16{,}67\),
\(\operatorname{MSE}=1100/3\approx366{,}67\),
\(\operatorname{RMSE}\approx19{,}15\).
Числа служат только для проверки понимания формул.
\end{keybox}

Дополнительно можно вычислить коэффициент детерминации
\begin{equation}
R^2=1-\frac{\sum_i(y_i-\widehat y_i)^2}
{\sum_i(y_i-\overline y)^2}.
\end{equation}
Здесь \(\overline y\) --- среднее по \textit{оцениваемой} части данных
в определении метрики, а не параметр предобработки.
Значение \(R^2\) может быть отрицательным; это не процент точных ответов.
При постоянной цели знаменатель равен нулю, и такая формула не определена
~\cite{metrics}.

\textbf{Функция потерь} задаёт то, что минимизирует обучение.
\textbf{Метрика} нужна для оценки и сравнения. При смене MSE на MAE
нельзя сравнивать сами значения \file{loss} как одну общую шкалу.
Во всех опытах сравнивайте одинаковые MAE и RMSE после возврата
к исходным ценам. Главным критерием выбора в этой работе назначается MAE
валидации; RMSE помогает увидеть крупные ошибки.

Для преобразования~\eqref{eq:target} справедливо
\(\operatorname{MAE}_y=s_y\operatorname{MAE}_z\) и
\(\operatorname{MSE}_y=s_y^2\operatorname{MSE}_z\).
Поэтому метрика в журнале обучения на стандартизованной цели и ошибка
в денежных единицах имеют разные численные значения.

\subsection{Обучение: градиент, эпоха и мини-выборка}
Обучение изменяет веса так, чтобы уменьшить среднюю потерю. Градиент
показывает, в какую сторону потеря растёт быстрее при малом изменении
параметров. Шаг в противоположную сторону помогает её уменьшить; величину
шага задаёт скорость обучения.
Для мини-выборки \(B_t\) шаг обычного градиентного метода имеет вид
\begin{equation}
\theta_{t+1}=\theta_t-\eta\nabla_\theta
\frac{1}{|B_t|}\sum_{i\in B_t}\ell(y_i,f_\theta(x_i)),
\end{equation}
где \(\eta>0\) --- скорость обучения. Производные вычисляются с
использованием цепного правила; последовательное распространение
производных от выхода сети к её входным слоям называют обратным
распространением ошибки.

\begin{keybox}{Один шаг обучения на простом примере}
Пусть \(\widehat y=wx+b\), \(x=2\), правильный ответ \(y=6\),
начальные \(w=1\), \(b=0\). Прогноз равен 2, а квадратичная потеря
\((\widehat y-y)^2=16\). Производные:
\[
\frac{\partial L}{\partial w}=2(\widehat y-y)x=-16,\qquad
\frac{\partial L}{\partial b}=2(\widehat y-y)=-8.
\]
При шаге \(0{,}01\) получаем \(w=1{,}16\), \(b=0{,}08\).
Новый прогноз равен \(2{,}4\), потеря --- \(12{,}96\): на этом шаге
она уменьшилась. В многослойной сети производные передаются через слои
по цепному правилу. Это учебный расчёт для одного объекта, не результат
обучения на домах.
\end{keybox}

\textbf{Итерация} в рассматриваемом режиме --- один шаг обновления весов
по мини-выборке. \textbf{Эпоха} --- проход по всей обучающей части.
Если в ней \(N\) строк и последний неполный пакет сохраняется, число
обновлений за эпоху равно \(\lceil N/b\rceil\), где \(b\) ---
\file{batch\_size}. При \(N=1000\), \(b=32\) это 32 обновления,
а за 10 эпох --- 320~\cite{training}.

Меньший пакет означает больше обновлений за эпоху и обычно более шумную
оценку градиента. Поэтому сравнение пакетов при одинаковом числе эпох
не является сравнением при одинаковом числе обновлений.
Слишком большая скорость обучения может привести к колебаниям или
расходимости; слишком малая --- к медленному улучшению.

В работе сравниваются SGD и Adam. SGD использует градиент текущего
пакета; Adam дополнительно накапливает оценки первого и второго моментов
градиента и адаптирует масштаб обновлений~\cite{sgd,adam}.
Одинаковая численная скорость обучения не означает одинаковые траектории
или одинаковую оптимальность настройки для двух алгоритмов.

\subsection{Переобучение и воспроизводимость}
\textbf{Регуляризация} ограничивает подгонку модели к особенностям
обучения. Например, штраф L2 добавляет к потере величину, растущую с
квадратами весов. Dropout во время обучения случайно отключает часть
выходов скрытого слоя, а при прогнозировании отключение не выполняется.
Эти приёмы меняют обучение и требуют проверки на валидации; в базовую
конфигурацию они не включены.

Переобучение можно заподозрить, если ошибка на обучении продолжает
уменьшаться, а на валидации перестаёт улучшаться или растёт.
По одной кривой нельзя автоматически установить причину: проверьте
подготовку данных, размер выборки, архитектуру и скорость обучения.

Дополнительный приём --- ранний останов: прекратить обучение после
заданного числа эпох без улучшения валидационной метрики и вернуть
лучшие сохранённые веса. В Keras это реализует
\file{EarlyStopping}~\cite{early}. В основных опытах с числом эпох
ранний останов выключен, чтобы сравнивать результаты после заданного
числа проходов. Его эффект исследуется отдельно.

Используйте одинаковое разбиение и исходный seed 40.
Начальные состояния генераторов задаются командой
\file{keras.utils.set\_random\_seed}.
В примере дополнительно включаются детерминированные операции, что уменьшает
дополнительную случайность вычислений~\cite{seed}.
Записывайте версии библиотек и устройство. Один seed не является
оценкой устойчивости и не гарантирует побитового совпадения на разных
программных и аппаратных платформах~\cite{determinism}.
